% ==============================================================================
% 90-appendices.tex — Приложения A, B, C
% ==============================================================================

\chapter{Макеты пользовательского интерфейса}
\label{app:ui}

Экраны приложения: главная страница, вход в~систему, регистрация, чат с~ассистентом, профиль бизнеса, управление подключением каналов, лента отзывов, история публикаций, история задач, настройки учётной записи.

\chapter{Модели данных}
\label{app:data-model}

Хранение данных: реляционная СУБД (пользователи, бизнесы, расписания, интеграции, подписки, журнал аудита, токены обновления, статусы синхронизации по платформам); документная СУБД (диалоги, сообщения, задачи агентов, отзывы, публикации). Инфологическая модель (ER-диаграмма)~--- в~пояснительной записке ВКР, глава~1.

\chapter{Дополнительная информация}
\label{app:extra}

Классификация целевых платформ по типу интеграции~--- таблица~\ref{tab:platform-class}.

\begin{table}[H]
  \caption{Платформы по типу интеграции}
  \label{tab:platform-class}
  \small
  \begin{tabularx}{\textwidth}{|p{3cm}|p{3.5cm}|X|}
    \hline
    \textbf{Платформа} & \textbf{Метод} & \textbf{Примечание} \\
    \hline
    VK & Программный интерфейс (официальный) & OAuth~2.0 \\
    Telegram & Интерфейс бота (официальный) & Токен бота \\
    Яндекс.Бизнес & Браузерная автоматизация & Нет публичного интерфейса \\
    2ГИС & Браузерная автоматизация & Нет публичного интерфейса \\
    Google Business & Программный интерфейс (официальный) & Для международных рынков \\
    \hline
  \end{tabularx}
\end{table}
